%% 自動生成（gen.py）
\bktocchap{9-0}{音を分ける、変化を予想する}{声を成分に分けると、何が見える？}{4}{3}{3}{8}{J}
\bktocpart{第1部}{フーリエ解析}{36}
\bktocchap{9-1}{波をたし合わせて形を作る：フーリエ級数}{sinをたしていくと、四角い波ができる}{4}{4}{4}{38}{C}
\bktocchap{9-2}{成分の取り出し方：直交性}{かけて積分すると、ほしい成分だけが残る}{4}{5}{5}{82}{C}
\bktocchap{9-3}{複素数で書くフーリエ級数}{sinとcosを、1つの式でまとめて書く}{4}{5}{5}{124}{C}
\bktocchap{9-4}{スペクトル：フーリエ変換}{くり返さない音も、周波数に分けられる}{7}{5}{5}{164}{C}
\bktocpart{第2部}{ラプラス変換}{212}
\bktocchap{9-5}{微分をかけ算に変える翻訳機：ラプラス変換}{微分方程式が、ふつうの方程式に変わる}{3}{5}{5}{214}{C}
\bktocchap{9-6}{変換表と性質}{よく使う関数の「翻訳表」を作ろう}{3}{5}{5}{248}{C}
\bktocchap{9-7}{ラプラス変換で微分方程式を解く}{部分分数に分けて、表で元にもどす}{3}{5}{5}{292}{C}
\bktocchap{9-8}{伝達関数：入れると出てくるを s で表す}{箱をつなぐと、伝達関数のかけ算になる}{3}{5}{5}{346}{C}
\bktocchap{9-9}{極と安定性：s の値でふるまいがわかる}{極の場所を見るだけで、落ち着くかどうかがわかる}{3}{5}{5}{388}{C}
\bktocpart{第3部}{行列と微分方程式}{442}
\bktocchap{9-10}{連立微分方程式を行列で：状態方程式}{たくさんの微分方程式を、1行で書く}{5}{5}{5}{444}{C}
\bktocchap{9-11}{行列の指数関数と固有値}{固有値が、落ち着くかどうかを決める}{5}{5}{5}{492}{C}
\bktocchap{9-12}{フーリエとラプラスの使いみち}{音楽の圧縮、機械の診断、エアコンの制御}{7}{3}{5}{538}{R}
\bktocchap{9-13}{工学の数学の地図}{7本の路線は、どこで合流した？}{all}{3}{5}{556}{E}
\par\vspace{2mm}
\bktocbacklist{\bkbi{こたえ}{582}\bkbi{\R{記号}{き|ごう}カード}{593}\bkbi{よく\R{使}{つか}う\R{定数}{てい|すう}・\R{単位}{たん|い}・\R{公式}{こう|しき}}{600}\bkbi{できるようになったことチェック\R{表}{ひょう}}{621}\bkbi{物理シリーズへ}{624}\bkbi{さくいん}{626}}
